% 供「二、问题分析」中问题三小节直接粘贴。

\paragraph{问题三。}
问题三的核心困难不是再列一套储能约束，而是在预测分四次到达时，决定\textbf{哪些普通购电承诺仍可改、改写后如何相对于 0:00 原计划付费}。由问题二可知，全日锁定的日前额度能够把“计划”和“执行”分开，但会放弃题目明确给予的 6:00、12:00、18:00 更新权，从而把本可在正常电价下完成的纠偏推迟为 5 倍紧急购电。因此首先冻结信息集：更新时点只用已发布的附件 3 光伏预报和日初因果负荷预报，已执行时段的承诺与 SOC 全部锁定；随后以非对称调整费 \(\phi_t(g^0,g^F)=p_tg^F+0.5p_t\lvert g^F-g^0\rvert\) 为目标中的计划项，对剩余时域建立线性规划；最后用“维持原承诺”与“允许调整”的预测成本差 \(VoI\) 记录是否实施新计划，并以从不调整、只开放单一时点的消融回答新预报值不值得用。输出是全年（及题设 2--12 月）的计划购电量、调整购电量、充放电和紧急购电，以及可追溯的成本分解。问题二若替换主预测器或年末库存，只修改承接句，不修改本问的结算与窗口定义。
